\VerdierPhysicalPageAnchor{0027}%
éléments de $\operatorname{Ext}_{\mathcal A}^{n}(X,Y)$, $n$-ième hyper-ext
des complexes $X$ et $Y$ (au sens de [1]). Lorsque $X$ et $Y$ sont des
complexes bornés (composants nuls sauf un nombre fini), les homomorphismes de
degré $n$ de $X$ dans $Y$ peuvent se définir par une construction qui
généralise la construction de Yoneda [3]. La composition des morphismes dans
$\mathbf D(\mathcal A)$ n'est autre que la composition de Yoneda.

\VerdierSemanticAnchor{intro-6-plan-chapter-three-spectral-objects}
Au paragraphe 4, nous étudions les objets spectraux usuels. Tout complexe de
$\mathcal A$ donne naissance à trois objets spectraux à valeurs dans
$\mathbf D(\mathcal A)$, qui permettent d'écrire de façon automatique les
suites spectrales essentielles. D'ailleurs toute suite
$\cdots\to X_n\to X_{n+1}\to\cdots$ de complexes et de morphismes de
complexes (où les morphismes de transition ne sont pas nécessairement des
monomorphismes) donne naissance à un objet spectral.

\VerdierSemanticAnchor{intro-7}
\noindent\textbf{7.}\quad
En résumé, il s'agit d'établir les fondements d'un formalisme, avec tous les
inconvénients qu'un tel travail comporte : enthymèmes et sorites,
démonstrations sans réelles difficultés mathématiques mais nécessitant
souvent des suites de vérifications parfois (resp. toujours) ennuyeuses. Ce
travail ne présente donc un intérêt que dans la mesure où ce formalisme, par
sa souplesse et sa généralité, permet la formulation et la démonstration de
\og vrais\fg{} théorèmes. C'est le souci d'énoncer et de démontrer des
théorèmes de dualité cohomologique qui a constitué la première motivation de
ce formalisme. On sait qu'il existe dans différents contextes des théories de
dualité cohomologique formellement très analogues :

\VerdierSemanticAnchor{intro-7-duality-theories}
\begin{enumerate}
\renewcommand{\labelenumi}{\arabic{enumi})}
\VerdierSemanticAnchor{intro-7-duality-1}
\item Théorie de la dualité dans la cohomologie des faisceaux cohérents sur
les schémas [4].

\VerdierSemanticAnchor{intro-7-duality-2}
\item Théorie de la dualité dans la cohomologie des faisceaux pour la
topologie étale des schémas [5], [6].

Ces théories sont dues à A. Grothendieck.

\VerdierSemanticAnchor{intro-7-duality-3}
\item Théorie de la dualité dans la cohomologie des modules galoisiens due à
J. Tate et J.-L. Verdier [7].

\VerdierSemanticAnchor{intro-7-duality-4}
\item Théorie de la dualité dans la cohomologie des faisceaux sur un espace
localement compact [9], [10].

\VerdierSemanticAnchor{intro-7-duality-5}
\item Théorie de la dualité dans la cohomologie des faisceaux analytiques
cohérents. Cette théorie reste encore essentiellement conjecturale.

\VerdierSemanticAnchor{intro-7-duality-6}
\item Théorie de la dualité pour les corps locaux, corps de nombres ou corps
de fonctions et ses relations avec la théorie du corps de classes. Cette
théorie est à développer. Les premiers pas sont dus à M. Artin [8].
\end{enumerate}

\VerdierSemanticAnchor{intro-7-adjunction}
Toutes ces théories ont pour point de départ (ou d'arrivée) une propriété
d'adjonction reliant certains foncteurs entre catégories dérivées. Le

